\section{State fibres and characteristic lifts}

A state $U_*=(u_*,v_*)\in\mathcal U$ is not, in general, a point of the wave
manifold.  For $P\in\mathcal W$, the canonical projections
\[
 \pi_-(P)=U_-,\qquad \pi_+(P)=U_+
\]
associate to it the inverse images
\[
 \mathcal F_-(U_*):=\pi_-^{-1}(U_*),\qquad
 \mathcal F_+(U_*):=\pi_+^{-1}(U_*).
\]
Generically these inverse images are curves in $\mathcal W$.  The notation
$\mathcal F_\pm$ is used here only for projection fibres; it does not introduce
a Hugoniot notation or an admissibility condition.

\subsection{Finite chart}
Let
\[
 L_{U_*}(z):=P(z)u_*+Q(z)v_*+R(z),\qquad
 D_r(z):=P(z)z+Q(z).
\]
For a point of $\mathcal F_\sigma(U_*)$, where $\sigma=-1$ denotes $\pi_-$
and $\sigma=+1$ denotes $\pi_+$,
\[
 U_*=\bar U+\frac{\sigma Y}{2}(z,1).
\]
Substitution in $P\bar u+Q\bar v+R=0$ gives, whenever $D_r\ne0$,
\[
 \boxed{Y=\frac{2L_{U_*}(z)}{\sigma D_r(z)}}.
\]
The midpoint follows from
\[
 \bar U=U_*-\frac{\sigma Y}{2}(z,1),
\]
and $\tau$ is recovered uniquely from
\[
 \bar U=U_0(z)+\tau(-Q(z),P(z)).
\]
Thus $z$ parametrizes a regular part of the projection fibre.  If $D_r=0$,
this parametrization can fail or become non-unique; this is recorded explicitly
rather than silently choosing a branch.

\subsection{Full reciprocal transition}
The reciprocal chart must also transform the jump coordinate.  Since
\[
 (z,1)=Z^{-1}(1,Z),\qquad Z=1/z,
\]
write the reciprocal jump coordinate as $\Upsilon$.  Equality of the physical
end states on the chart overlap requires
\[
 \boxed{\eta=\frac{\tau}{Z^2}},\qquad
 \boxed{\Upsilon=\frac{Y}{Z}}.
\]
Hence the full reciprocal wave coordinates are $(\eta,\Upsilon,Z)$, not
$(\eta,Y,Z)$ with an unchanged numerical $Y$.  On the characteristic both
jump coordinates vanish, so the earlier rarefaction calculations on $Y=0$
are unaffected by this clarification.

With
\[
 \widetilde L_{U_*}(Z)=\widetilde P(Z)u_*+\widetilde Q(Z)v_*+\widetilde R(Z),
 \qquad
 \widetilde D_r(Z)=\widetilde P(Z)+Z\widetilde Q(Z),
\]
the regular reciprocal fibre satisfies
\[
 \boxed{\Upsilon=\frac{2\widetilde L_{U_*}(Z)}{\sigma\widetilde D_r(Z)}}.
\]
At $Z=0$ the denominator can vanish.  If the numerator vanishes as well, the
fixed-state constraint may cease to determine $\Upsilon$ uniquely.  This is a
genuine projective degeneracy of the fibre parametrization and is not replaced
by a fabricated finite-$z$ point.

\subsection{Characteristic lifts of a state}
The intersections of either projection fibre with the characteristic satisfy
$Y=0$ (or $\Upsilon=0$), hence
\[
 \mathcal F_-(U_*)\cap\mathcal C
 =\mathcal F_+(U_*)\cap\mathcal C.
\]
In the finite chart their directions are the real roots of
\[
 \boxed{L_{U_*}(z)=A_2z^2+A_1z+A_0=0},
\]
where
\begin{align*}
 A_2&=2b_{20}u_*+b_{11}v_*+b_{10},\\
 A_1&=(b_{11}-2a_{20})u_*+(2b_{02}-a_{11})v_*+(b_{01}-a_{10}),\\
 A_0&=-a_{11}u_*-2a_{02}v_*-a_{01}.
\end{align*}
These are precisely the projective characteristic directions at $U_*$.  If
$A_2=0$, the projective root at infinity is represented in the reciprocal chart
by $Z=0$.  Therefore a state may possess more than one characteristic lift;
it must not be identified with a unique point of $\mathcal C$ or $\mathcal W$.
